% ==============================================================================
% Chapter 2: Valuation of Financial Instruments
% ==============================================================================
\chapter{Valuation of Financial Instruments}

% ==============================================================================
\section{Foundations of Valuation and Yield Concepts}
\lecturedate{07/09/2026}
% ==============================================================================

\begin{definition}{Foundational Valuation Concepts}{val_concepts}
In corporate finance and financial mathematics, asset valuation depends fundamentally on the perspective and purpose of evaluation:
\begin{enumerate}[leftmargin=*]
    \item \textbf{Book Value:} The accounting net worth of an asset or firm recorded on the balance sheet (historic cost minus accumulated depreciation).
    \item \textbf{Market Value:} The actual price at which an asset or financial security can be bought or sold in an active, competitive market.
    \item \textbf{Going Concern Value:} The economic value a prospective purchaser is prepared to pay for an operating business with established operations and recurring earnings power.
    \item \textbf{Liquidating Value (Breakup Value):} The net residual sum realized if all assets of the enterprise are sold piecemeal and all external liabilities are settled.
    \item \textbf{Capitalised Value:} The present value of all anticipated future cash inflows generated by a fixed asset or project, discounted at the firm's required cost of capital.
\end{enumerate}
\end{definition}

\subsection{Yield Fundamentals}

\begin{formulabox}[Market Yield and Current Yield]
The \textbf{Current Yield} measures the annual income return generated by a fixed-income security or preference share relative to its current market price:
\begin{equation}
    \text{Current Yield} = \frac{\text{Annual Dividend / Coupon}}{\text{Current Market Price}} \times 100
\end{equation}
The \textbf{Dividend Yield} ($dY$) of an equity or preference share is defined as:
\begin{equation}
    dY = \frac{\DPS}{P_0} \times 100
\end{equation}
where $\DPS$ is the annual dividend per share and $P_0$ is the current market price per share.
\end{formulabox}

\begin{example}{Current Yield of Preference Shares}{current_yield_pref}
An $8\%$ preference share of ABC Limited with face value $\inr 100$ is selling in the market for $\inr 64$. Determine the current yield of the preference share.
\tcblower
\textbf{Formula:}
\begin{equation*}
    \text{Current Yield} = \frac{\text{Annual Dividend}}{P_0} \times 100
\end{equation*}
\textbf{Substitution:}
\begin{equation*}
    \text{Current Yield} = \frac{8}{64} \times 100 = 0.125 \times 100 = 12.5\%
\end{equation*}
\textbf{Value:} $12.5\%$
\end{example}

\begin{example}{Expected Return on XYZ Limited Preference Shares}{exp_ret_xyz}
The preference shares of XYZ Ltd are selling for $\inr 50$ per share and pay an annual dividend of $\inr 3$. What is the expected rate of return if the security is purchased at the market price?
\tcblower
\textbf{Formula:}
\begin{equation*}
    \text{Expected Return} = \frac{\DPS}{P_0} \times 100
\end{equation*}
\textbf{Substitution:}
\begin{equation*}
    \text{Expected Return} = \frac{3}{50} \times 100 = 6\%
\end{equation*}
\textbf{Value:} $6\%$
\end{example}

\begin{example}{Valuation and Investment Decision (Bhuvanesh / PQR Ltd)}{inv_decision_bhuvanesh}
Bhuvanesh owns 200 preference shares of PQR Ltd, currently selling at $\inr 40$ per share and paying an annual dividend of $\inr 7$ per share.
\begin{enumerate}[leftmargin=*, noitemsep]
    \item What is the expected return (dividend yield)?
    \item If Bhuvanesh requires a $12\%$ return, should he sell or buy more preference shares?
\end{enumerate}
\tcblower
\textbf{1. Expected Return:}
\begin{equation*}
    \text{Expected Return} = \frac{7}{40} \times 100 = 17.5\%
\end{equation*}
\textbf{2. Intrinsic Valuation and Decision Rule:}
\begin{equation*}
    P^* = \frac{D}{K} = \frac{7}{0.12} = \inr 58.33
\end{equation*}
Since the expected rate of return ($17.5\%$) exceeds the required rate of return ($12\%$) and the market price ($P_0 = \inr 40$) is below intrinsic worth ($P^* = \inr 58.33$), the security is undervalued. Hence, Bhuvanesh should \textbf{buy more preference shares}.
\end{example}

% ==============================================================================
\section{Yield to Maturity, Holding Period Return, and Equity Valuation}
\lecturedate{08/09/2026}
% ==============================================================================

\begin{definition}{Yield Measures and Holding Period Return}{yield_measures}
\paragraph{Yield to Maturity (YTM / Redemption Yield)}
The overall annualized rate of return earned by an investor who holds a fixed-income bond or redeemable preference share to its maturity date. YTM reflects:
\begin{enumerate}[noitemsep, leftmargin=*]
    \item Recurring annual coupon/dividend cash flows.
    \item Periodic capital appreciation or depreciation over the holding span (difference between face value and purchase price).
\end{enumerate}

\paragraph{Holding Period Return (HPR)}
The total return achieved over a planned holding duration:
\begin{equation}
    \HPR = \frac{D_t + \Delta P}{P_1} = \frac{D_t + (P_t - P_1)}{P_1}
\end{equation}
where $D_t$ is the dividend received during the period, $P_1$ is the purchase price, $P_t$ is the terminal selling price, and $\Delta P = P_t - P_1$ represents capital gain/loss.

\paragraph{Earnings Yield}
The proportion of accounting net earnings generated per share relative to market price:
\begin{equation}
    \text{Earnings Yield} = \frac{\EPS}{P_0}
\end{equation}
\begin{itemize}[noitemsep]
    \item \textbf{Gross Yield:} Used by corporate treasurers to estimate the effective cost of equity capital.
    \item \textbf{Net Yield:} Used by equity investors to evaluate the true underlying investment return after taxes.
\end{itemize}
\end{definition}

\subsection{Approaches to Equity Share Valuation}

\begin{formulabox}[Methods of Equity Valuation]
Equity valuation methodologies fall into four structured frameworks:
\begin{enumerate}[noitemsep, leftmargin=*]
    \item \textbf{Accounting-Based Models:} Book Value (balance sheet net worth) and Liquidation Value.
    \item \textbf{Dividend Discount Models (DDM):} Present value of all projected future dividend streams.
    \item \textbf{Discounted Cash Flow (DCF) Models:} Operating cash flows (FCFF) and Free Cash Flows to Equity (FCFE).
    \item \textbf{Relative Valuation Metrics:} Price-to-Earnings ($P/E$), Price-to-Book ($P/B$), and Price-to-Sales ($P/S$).
\end{enumerate}
\end{formulabox}

\subsection{The General Dividend Discount Model (DDM)}

\begin{definition}{General DDM Assumptions}{ddm_assumptions}
Under fundamental dividend discount theory:
\begin{enumerate}[noitemsep, leftmargin=*]
    \item Dividends are paid annually at discrete intervals.
    \item The first dividend $D_1$ is received exactly one year after acquisition.
    \item Equity liquidation occurs at year-end on ex-dividend terms.
\end{enumerate}
\end{definition}

\begin{formulabox}[Finite and Infinite DDM Formulations]
For a \textbf{Finite Holding Period ($n$ years)} with terminal market price $P_n$:
\begin{equation}
    P_0 = \sum_{t=1}^{n} \frac{D_t}{(1 + K_e)^t} + \frac{P_n}{(1 + K_e)^n}
\end{equation}
Specifically, for a \textbf{Two-Year Holding Horizon ($n = 2$)}:
\begin{equation}
    P_0 = \frac{D_1}{1 + K_e} + \frac{D_2 + P_2}{(1 + K_e)^2}
\end{equation}
For an \textbf{Infinite Holding Horizon ($n \to \infty$)}, share price equals the discounted sum of all future perpetual dividends:
\begin{equation}
    P_0 = \sum_{t=1}^{\infty} \frac{D_t}{(1 + K_e)^t}
\end{equation}
where $P_0$ is intrinsic share price, $D_t$ is expected dividend at year $t$, and $K_e$ is the required equity return.
\end{formulabox}

% ==============================================================================
\section{Dividend Valuation Models: Constant Dividend and Gordon Basics}
\lecturedate{10/09/2026}
% ==============================================================================

\subsection{Model 1: Zero Growth (Constant Perpetual Dividend)}

\begin{formulabox}[Zero Growth Valuation Formula]
When a firm pays an invariant perpetual dividend ($D_t = D$ for all $t$):
\begin{equation}
    P_0 = \frac{D}{K_e}
\end{equation}
\end{formulabox}

\begin{example}{Zero-Growth Equity Valuation}{zero_growth_val}
A firm pays an annual dividend of $20\%$ on equity shares of face value $\inr 100$ each. Find the intrinsic value of the share given that dividends remain constant perpetually and the required rate of return is $15\%$.
\tcblower
\textbf{Formula:}
\begin{equation*}
    P_0 = \frac{D}{K_e} = \frac{\text{Face Value} \times \text{Dividend Rate}}{K_e}
\end{equation*}
\textbf{Substitution:}
\begin{equation*}
    D = 100 \times 20\% = \inr 20, \quad K_e = 0.15
\end{equation*}
\begin{equation*}
    P_0 = \frac{20}{0.15} = \inr 133.33
\end{equation*}
\textbf{Value:} $\inr 133.33$
\end{example}

\subsection{Model 2: Constant Growth Dividend Model}

\begin{formulabox}[Gordon Constant Growth Formula]
Assuming dividends expand at a steady compound annual growth rate $g$ ($g < K_e$):
\begin{equation}
    P_0 = \frac{D_1}{K_e - g} = \frac{D_0 (1 + g)}{K_e - g}
\end{equation}
\end{formulabox}

\begin{example}{Constant Growth Share Valuation}{const_growth_val}
A share of face value $\inr 100$ pays an expected dividend of $12\%$ at the end of Year 1 ($D_1$). Dividend growth is estimated at $3\%$ per annum ($g$). If the investor's required rate of return is $16\%$ ($K_e$), calculate the intrinsic equity value.
\tcblower
\textbf{Formula:}
\begin{equation*}
    P_0 = \frac{D_1}{K_e - g}
\end{equation*}
\textbf{Substitution:}
\begin{equation*}
    D_1 = 100 \times 12\% = \inr 12, \quad K_e = 0.16, \quad g = 0.03
\end{equation*}
\begin{equation*}
    P_0 = \frac{12}{0.16 - 0.03} = \frac{12}{0.13} = \inr 92.31
\end{equation*}
\textbf{Value:} $\inr 92.31$
\end{example}

\subsection{Specific Dividend Valuation Models: Gordon's Model Fundamentals}

\begin{definition}{Gordon Model Assumptions and Growth Formulation}{gordon_foundations}
Under Myron J. Gordon's dividend capitalization framework:
\begin{enumerate}[noitemsep, leftmargin=*]
    \item Earnings are either distributed as cash dividends or retained for internal reinvestment.
    \item The retention ratio $b$ and rate of return on reinvested capital $r$ remain strictly constant.
    \item Growth rate in future earnings and dividends is endogenous:
    \begin{equation}
        g = b \times r
    \end{equation}
    \item The firm is an all-equity enterprise with infinite life, and $K_e > br$.
\end{enumerate}
\end{definition}

\begin{formulabox}[Gordon's Valuation Formula]
\begin{equation}
    P_0 = \frac{\EPS_1 (1 - b)}{K_e - b \cdot r} = \frac{D_1}{K_e - g}
\end{equation}
where $\EPS_1$ is expected earnings per share, $b$ is the retention ratio, $(1 - b)$ is the dividend payout ratio, $K_e$ is the cost of equity, and $r$ is the internal rate of return.
\end{formulabox}

\begin{example}{Earnings Reinvestment vs. Full Payout}{earnings_reinvest_gordon}
A firm generates $\EPS_1 = \inr 9$ and faces a cost of capital $K_e = 15\%$. Compare the expected dividend $D_1$ under:
\begin{enumerate}[noitemsep, leftmargin=*]
    \item Case 1: Retention ratio $b = 0$ (entire earnings distributed as dividend).
    \item Case 2: Payout ratio $1 - b = 60\%$ (retention ratio $b = 40\%$).
\end{enumerate}
\tcblower
\textbf{Case 1 ($b = 0$):}
\begin{equation*}
    D_1 = \EPS_1 \times (1 - b) = 9 \times (1 - 0) = \inr 9.00
\end{equation*}
\textbf{Case 2 ($b = 0.40$):}
\begin{equation*}
    D_1 = \EPS_1 \times (1 - b) = 9 \times 0.60 = \inr 5.40
\end{equation*}
\textbf{Value:} Case 1: $\inr 9.00$ \quad | \quad Case 2: $\inr 5.40$
\end{example}

% ==============================================================================
\section{Gordon's Model: Growth, Normal, and Declining Firm Dynamics}
\lecturedate{15/09/2026}
% ==============================================================================

\begin{formulabox}[Firm Classification and Payout Policy Implications]
Gordon's model establishes three distinct economic regimes based on the relationship between return on investment ($r$) and capitalization rate ($K_e$):
\begin{enumerate}[leftmargin=*, noitemsep]
    \item \textbf{Growth Firm ($r > K_e$):} The firm earns more on retained capital than equity holders could earn externally. The optimal dividend payout ratio is \textbf{NIL ($0\%$)}, and retention should be maximized ($b = 100\%$).
    \item \textbf{Normal Firm ($r = K_e$):} Internal returns match investors' required returns. The dividend payout ratio is \textbf{irrelevant}; market price remains identical across all payout levels.
    \item \textbf{Declining Firm ($r < K_e$):} Internal returns are lower than the opportunity cost of capital. The optimal dividend payout ratio is \textbf{$100\%$}, and retention should be zero ($b = 0\%$).
\end{enumerate}
\end{formulabox}

\begin{example}{Gordon's Model for a Growth Firm (XY Ltd)}{gordon_growth_xy}
XY Ltd is a growth firm with $\EPS = \inr 14$, capitalization rate $K_e = 15\%$, and internal rate of return $r = 20\%$. Determine the market price per share under Gordon's model for retention ratios of: (a) $40\%$, (b) $60\%$, (c) $20\%$.
\tcblower
\textbf{General Formula:}
\begin{equation*}
    P_0 = \frac{\EPS \times (1 - b)}{K_e - b \cdot r} = \frac{D}{K_e - g}
\end{equation*}
\textbf{(a) Retention $b = 40\%$ (Payout $= 60\%$):}
\begin{align*}
    D &= 14 \times 0.60 = \inr 8.40, \quad g = 0.40 \times 0.20 = 0.08 \\
    P_0 &= \frac{8.40}{0.15 - 0.08} = \frac{8.40}{0.07} = \inr 120.00
\end{align*}
\textbf{(b) Retention $b = 60\%$ (Payout $= 40\%$):}
\begin{align*}
    D &= 14 \times 0.40 = \inr 5.60, \quad g = 0.60 \times 0.20 = 0.12 \\
    P_0 &= \frac{5.60}{0.15 - 0.12} = \frac{5.60}{0.03} = \inr 186.67
\end{align*}
\textbf{(c) Retention $b = 20\%$ (Payout $= 80\%$):}
\begin{align*}
    D &= 14 \times 0.80 = \inr 11.20, \quad g = 0.20 \times 0.20 = 0.04 \\
    P_0 &= \frac{11.20}{0.15 - 0.04} = \frac{11.20}{0.11} = \inr 101.82
\end{align*}
\textbf{Conclusion:} Because $r > K_e$ ($20\% > 15\%$), market price per share increases as the retention ratio increases. The growth firm maximizes shareholder wealth by retaining maximum earnings.
\end{example}

\begin{example}{Gordon's Model for a Normal Firm (MN Ltd)}{gordon_normal_mn}
MN Ltd is a normal firm with $\EPS = \inr 45$, cost of capital $K_e = 18\%$, and return on investment $r = 18\%$. Ascertain the market value per share using Gordon's model if dividend payout is: (a) $30\%$, (b) $60\%$, (c) $90\%$.
\tcblower
\textbf{General Formula:}
\begin{equation*}
    P_0 = \frac{\EPS \times (1 - b)}{K_e - b \cdot r}
\end{equation*}
\textbf{(a) Payout $= 30\%$ (Retention $b = 70\%$):}
\begin{equation*}
    D = 45 \times 0.30 = \inr 13.50, \quad g = 0.70 \times 0.18 = 0.126
\end{equation*}
\begin{equation*}
    P_0 = \frac{13.50}{0.18 - 0.126} = \frac{13.50}{0.054} = \inr 250.00
\end{equation*}
\textbf{(b) Payout $= 60\%$ (Retention $b = 40\%$):}
\begin{equation*}
    D = 45 \times 0.60 = \inr 27.00, \quad g = 0.40 \times 0.18 = 0.072
\end{equation*}
\begin{equation*}
    P_0 = \frac{27.00}{0.18 - 0.072} = \frac{27.00}{0.108} = \inr 250.00
\end{equation*}
\textbf{(c) Payout $= 90\%$ (Retention $b = 10\%$):}
\begin{equation*}
    D = 45 \times 0.90 = \inr 40.50, \quad g = 0.10 \times 0.18 = 0.018
\end{equation*}
\begin{equation*}
    P_0 = \frac{40.50}{0.18 - 0.018} = \frac{40.50}{0.162} = \inr 250.00
\end{equation*}
\textbf{Conclusion:} For a normal firm ($r = K_e = 18\%$), the share price is invariant to payout choice ($P_0 = \inr 250.00$). Dividend policy is completely neutral.
\end{example}

\begin{example}{Gordon's Model for a Declining Firm (ABC Ltd)}{gordon_declining_abc}
ABC Ltd is a declining firm with $\EPS = \inr 35$, capitalization rate $K_e = 15\%$, and productivity of retained earnings $r = 10\%$. Using Gordon's model, determine the market price per share if payout is: (a) $20\%$, (b) $40\%$, (c) $70\%$.
\tcblower
\textbf{General Formula:}
\begin{equation*}
    P_0 = \frac{\EPS \times (1 - b)}{K_e - b \cdot r}
\end{equation*}
\textbf{(a) Payout $= 20\%$ (Retention $b = 80\%$):}
\begin{equation*}
    D = 35 \times 0.20 = \inr 7.00, \quad g = 0.80 \times 0.10 = 0.08
\end{equation*}
\begin{equation*}
    P_0 = \frac{7.00}{0.15 - 0.08} = \frac{7.00}{0.07} = \inr 100.00
\end{equation*}
\textbf{(b) Payout $= 40\%$ (Retention $b = 60\%$):}
\begin{equation*}
    D = 35 \times 0.40 = \inr 14.00, \quad g = 0.60 \times 0.10 = 0.06
\end{equation*}
\begin{equation*}
    P_0 = \frac{14.00}{0.15 - 0.06} = \frac{14.00}{0.09} = \inr 155.56
\end{equation*}
\textbf{(c) Payout $= 70\%$ (Retention $b = 30\%$):}
\begin{equation*}
    D = 35 \times 0.70 = \inr 24.50, \quad g = 0.30 \times 0.10 = 0.03
\end{equation*}
\begin{equation*}
    P_0 = \frac{24.50}{0.15 - 0.03} = \frac{24.50}{0.12} = \inr 204.17
\end{equation*}
\textbf{Conclusion:} Because $r < K_e$ ($10\% < 15\%$), market price per share increases monotonically as the dividend payout ratio increases. The optimal policy for a declining firm is $100\%$ payout.
\end{example}

% ==============================================================================
\section{Walter's Dividend Valuation Model and Policy Decisions}
\lecturedate{16/09/2026}
% ==============================================================================

\subsection{Theories of Dividend Policy}

\begin{definition}{Dividend Relevance vs. Irrelevance}{theories_dividend}
\begin{itemize}[leftmargin=*, noitemsep]
    \item \textbf{Theories of Relevance:} Maintain that dividend payouts actively dictate share valuation and firm wealth. Key formulations: \emph{Gordon's Model} and \emph{Walter's Model}.
    \item \textbf{Theories of Irrelevance:} Posit that dividend distributions do not affect shareholder wealth or the cost of capital in perfect markets. Key formulation: \emph{Modigliani-Miller (MM) Hypothesis}.
\end{itemize}
\end{definition}

\subsection{Prof. E. James Walter's Model}

\begin{definition}{Walter's Model Foundations and Assumptions}{walter_assumptions}
Prof. E. James Walter argued that dividend policy is an inseparable investment decision. The relationship between the firm's internal rate of return ($r$) and its capitalization rate ($k$) determines value:
\begin{enumerate}[noitemsep, leftmargin=*]
    \item \textbf{Internal Equity Financing:} Retained earnings represent the sole financing source; no new equity or debt is issued.
    \item \textbf{Constant Returns:} Return on investment $r$ and cost of capital $k$ remain constant over time.
    \item \textbf{Infinite Horizon:} The firm has an infinite economic life.
    \item \textbf{Constant Earnings and Dividends:} $\EPS$ and $\DPS$ are invariant in assessing baseline values.
\end{enumerate}
\end{definition}

\begin{formulabox}[Walter's Valuation Formula]
The theoretical market price per share ($P$) under Walter's model is given by:
\begin{equation}
    P = \frac{D + \frac{r}{k}(E - D)}{k}
\end{equation}
where:
\begin{itemize}[noitemsep]
    \item $P$ = Theoretical market price per share ($\MPS$)
    \item $D$ = Dividend per share ($\DPS$)
    \item $E$ = Earnings per share ($\EPS$)
    \item $r$ = Firm's internal rate of return on investment
    \item $k$ = Cost of capital / capitalization rate
    \item $(E - D)$ = Retained earnings per share
    \item $\frac{r}{k}(E - D)$ = Capitalized market worth of retained earnings
\end{itemize}
Total Value of the Firm ($V$) with $N$ outstanding equity shares is:
\begin{equation}
    V = N \times P
\end{equation}
\end{formulabox}

\begin{example}{Walter's Model for a Growth Firm (Nirmal Ltd)}{walter_growth_nirmal}
Nirmal Ltd has a cost of capital $k = 10\%$ and rate of return on investment $r = 18\%$. The company has $500{,}000$ outstanding equity shares and an $\EPS = \inr 20$. Compute the market price per share and total value of the firm under Walter's model for: (a) No retention ($0\%$), (b) $40\%$ retention, (c) $80\%$ retention.
\tcblower
\textbf{Formula:}
\begin{equation*}
    P = \frac{D + \frac{r}{k}(E - D)}{k}, \quad V = N \times P
\end{equation*}
Here $E = \inr 20$, $r = 0.18$, $k = 0.10$, so $\frac{r}{k} = \frac{0.18}{0.10} = 1.80$.

\textbf{(a) Retention $0\%$ (Payout $100\%, D = 20$):}
\begin{align*}
    P &= \frac{20 + 1.80(20 - 20)}{0.10} = \frac{20 + 0}{0.10} = \inr 200.00 \\
    V &= 500{,}000 \times 200 = \inr 10{,}00{,}00{,}000 \quad (10\text{ Crore})
\end{align*}

\textbf{(b) Retention $40\%$ (Payout $60\%, D = 20 \times 0.60 = 12$):}
\begin{align*}
    P &= \frac{12 + 1.80(20 - 12)}{0.10} = \frac{12 + 14.40}{0.10} = \frac{26.40}{0.10} = \inr 264.00 \\
    V &= 500{,}000 \times 264 = \inr 13{,}20{,}00{,}000 \quad (13.2\text{ Crore})
\end{align*}

\textbf{(c) Retention $80\%$ (Payout $20\%, D = 20 \times 0.20 = 4$):}
\begin{align*}
    P &= \frac{4 + 1.80(20 - 4)}{0.10} = \frac{4 + 28.80}{0.10} = \frac{32.80}{0.10} = \inr 328.00 \\
    V &= 500{,}000 \times 328 = \inr 16{,}40{,}00{,}000 \quad (16.4\text{ Crore})
\end{align*}
\textbf{Conclusion:} Because $r > k$ ($18\% > 10\%$), Nirmal Ltd is a growth firm. As dividend payout increases, firm value diminishes. The optimal payout ratio is $0\%$ (full retention).
\end{example}

\begin{example}{Walter's Model for a Normal Firm}{walter_normal_firm}
The $\EPS$ of a company is $\inr 15$, its rate of capitalization $k = 12\%$, and return on investment $r = 12\%$. Compute the market value per share under Walter's model for payout ratios of: (a) $20\%$, (b) $50\%$, (c) $70\%$. Which is the optimum payout ratio?
\tcblower
\textbf{Formula:}
\begin{equation*}
    P = \frac{D + \frac{r}{k}(E - D)}{k}
\end{equation*}
Since $r = k = 12\%$, the ratio $\frac{r}{k} = \frac{0.12}{0.12} = 1.00$. Hence:
\begin{equation*}
    P = \frac{D + 1(15 - D)}{0.12} = \frac{15}{0.12} = \inr 125.00
\end{equation*}
\textbf{(a) Payout $20\%$ ($D = 15 \times 0.20 = \inr 3.00$):} $P = \inr 125.00$ \\
\textbf{(b) Payout $50\%$ ($D = 15 \times 0.50 = \inr 7.50$):} $P = \inr 125.00$ \\
\textbf{(c) Payout $70\%$ ($D = 15 \times 0.70 = \inr 10.50$):} $P = \inr 125.00$

\textbf{Conclusion:} Because $r = k$, share price remains constant at $\inr 125.00$ across all payout ratios. Dividend policy is neutral; there is no unique optimum payout ratio.
\end{example}

\begin{example}{Walter's Model for a Declining Firm}{walter_declining_firm}
The $\EPS$ of a firm is $\inr 12$, its capitalization rate is $k = 15\%$, and rate of return on investment is $r = 9\%$. Compute the market price per share using Walter's model for dividend payout ratios of: (a) $25\%$, (b) $50\%$, (c) $100\%$. Which is the ideal payout ratio?
\tcblower
\textbf{Formula:}
\begin{equation*}
    P = \frac{D + \frac{r}{k}(E - D)}{k}
\end{equation*}
Given $E = \inr 12$, $k = 0.15$, $r = 0.09$, so $\frac{r}{k} = \frac{0.09}{0.15} = 0.60$.

\textbf{(a) Payout $25\%$ ($D = 12 \times 0.25 = \inr 3.00$):}
\begin{equation*}
    P = \frac{3 + 0.60(12 - 3)}{0.15} = \frac{3 + 5.40}{0.15} = \frac{8.40}{0.15} = \inr 56.00
\end{equation*}
\textbf{(b) Payout $50\%$ ($D = 12 \times 0.50 = \inr 6.00$):}
\begin{equation*}
    P = \frac{6 + 0.60(12 - 6)}{0.15} = \frac{6 + 3.60}{0.15} = \frac{9.60}{0.15} = \inr 64.00
\end{equation*}
\textbf{(c) Payout $100\%$ ($D = 12 \times 1.00 = \inr 12.00$):}
\begin{equation*}
    P = \frac{12 + 0.60(12 - 12)}{0.15} = \frac{12 + 0}{0.15} = \inr 80.00
\end{equation*}
\textbf{Conclusion:} Because $r < k$ ($9\% < 15\%$), the company is a declining firm. Market price per share is maximized when all earnings are distributed. The ideal payout ratio is \textbf{$100\%$}.
\end{example}